(1) 取竖直向下为正，牛顿第二定律是 \(m\dot v=mg-kv^2\)，且 \(v(0)=0\)。定义终端速度 \(v_T=\sqrt{mg/k}\) 和时间尺度 \(\tau=v_T/g=\sqrt{m/(kg)}\)，则 \(\dot v=g[1-(v/v_T)^2]\)，适用于未发生方向反转的下落过程。
(2) 分离变量 \(\int_0^{v/v_T}du/(1-u^2)=gt/v_T\)，从而 \(\operatorname{artanh}(v/v_T)=t/\tau\)。故 \(v(t)=v_T\tanh(t/\tau)\)，位移 \(x(t)-x(0)=v_T\tau\ln\cosh(t/\tau)=(v_T^2/g)\ln\cosh(gt/v_T)\)。
(3) 当 \(\dot v=0\)，有 \(mg=kv_T^2\)，故 \(v_T=\sqrt{mg/k}\)。阻力瞬时耗散功率 \(P_{\rm diss}=kv^3\ge0\)；将动力学方程乘以 \(v\) 得 \(mgv=m v\dot v+kv^3=d(mv^2/2)/dt+P_{\rm diss}\)，证明重力做功被分成动能增长与热耗散。
(4) 在 \(t\ll\tau\)，\(\tanh(t/\tau)\sim t/\tau\)，故 \(v\sim gt\)、\(x-x_0\sim gt^2/2\)。\(t\gg\tau\) 时 \(v\to v_T\)、\(x\simeq v_Tt-(v_T^2/g)\ln2\)。固定时间令 \(k\to0\)，有 \(v_T\to\infty\) 但 \(t/\tau\to0\)，结合展开仍恢复匀加速下落；不能将极限理解为无限大的实际速度。
